% Part: computability
% Chapter: computability-theory
% Section: normal-form

\documentclass[../../../include/open-logic-section]{subfiles}

\begin{document}

\olfileid{cmp}{thy}{nfm}
\olsection{प्रमाण रूप प्रमेय}

संगणनक्षम फलनांच्या व्याख्यांचे वर्णन करता येते आणि एखाद्या आदानावर
त्या फलनाच्या संगणनाची पूर्ण नोंद असल्याचे सांगणारे कथित वर्णन खरे आहे
का हे तपासता येते, असे समजा. मग संगणनाच्या प्रतिमानावर अवलंबून न राहता,
कोणत्याही संगणनक्षम फलन~$f$ चे कोणत्याही आदान~$x$ वरील मूल्य पुढीलप्रमाणे
ठरवता येईल, असे वाटते:
\begin{enumerate}
  \item संगणनांच्या नोंदींच्या सर्व शक्य वर्णनांमध्ये शोध घेणे.
  \item एखादी नोंद~$f(x)$ च्या संगणनाची नोंद आहे का हे तपासणे.
  \item योग्य नोंद सापडल्यास तिच्यातून~$f(x)$ चे मूल्य काढणे.
\end{enumerate}
हे खरोखरच शक्य आहे, हाच क्लिनी यांच्या प्रमाण रूप प्रमेयाचा आशय आहे.

\begin{thm}[क्लिनीचे प्रमाण रूप प्रमेय]
\ollabel{thm:normal-form}
असा आदिम पुनरावर्ती संबंध~$T(e, x, s)$ आणि असे आदिम पुनरावर्ती
फलन~$U(s)$ असते की पुढील गुणधर्म पूर्ण होतो: $f$ हे कोणतेही
आंशिक संगणनक्षम फलन असेल, तर एखाद्या~$e$ साठी
\[
f(x) \simeq U(\umin{s}{T(e, x, s)})
\]
हे सर्व~$x$ साठी असते.
\end{thm}

\begin{proof}[सिद्धतेचे रूपरेषात्मक वर्णन]
संगणनाच्या कोणत्याही प्रतिमानासाठी संगणनक्षम फलन~$f$ चे वर्णन
काटेकोरपणे परिभाषित करून ते वर्णन नैसर्गिक संख्या~$e$ वापरून
संकेतबद्ध करता येते. त्याचप्रमाणे, निर्देशांक~$e$ असलेल्या फलनाचे
आदान~$x$ वरील संगणन कसे घडले याची नोंद करणाऱ्या ``संगणनक्रमिके''ची
कल्पना काटेकोरपणे परिभाषित करता येते. अशा संगणनक्रमिकेला देखील
संख्या~$s$ ने संकेतबद्ध करता येते. हे अशा प्रकारे करता येते की
\begin{enumerate}
  \item संख्या~$s$ ही निर्देशांक~$e$ असलेल्या फलनाच्या आदान~$x$
  वरील संगणनक्रमिकेचा संकेतांक असेल तेव्हाच लागू होणारा संबंध
  $T(e, x, s)$, आणि
  \item $s$ ने संकेतबद्ध केलेल्या संगणनक्रमिकेला त्या संगणनाचा अंतिम
  परिणाम देणारे फलन $U(s)$
\end{enumerate}
ही दोन्ही संगणनक्षम असतात. किंबहुना, संबंध~$T$ आणि फलन~$U$
आदिम पुनरावर्ती असतात.
\end{proof}

\begin{explain}
प्रमाण रूप प्रमेयाची काटेकोर सिद्धता द्यायची असेल, तर संगणनाचे एक
प्रतिमान ठरवून संगणनक्षम फलनांच्या वर्णनांचे आणि संगणनक्रमिकांचे
सांकेतिकरण तपशीलवार करावे लागेल; तसेच संबंध~$T$ आणि फलन~$U$
आदिम पुनरावर्ती आहेत हे पडताळावे लागेल. बहुतेक उपयोगांसाठी $T$
आणि~$U$ संगणनक्षम आहेत व $U$ सर्वत्र परिभाषित आहे, एवढे पुरेसे असते.

प्रमाण रूप प्रमेयाची सिद्धता एका संक्षिप्त सूत्रवाक्यात लक्षात ठेवणे
कदाचित उत्तम: $\umin{s}{T(e, x, s)}$ हे निर्देशांक~$e$ असलेल्या
फलनाच्या आदान~$x$ वरील संगणनक्रमिकेचा शोध घेते; अशी क्रमिका
सापडल्यास $U$ त्या संगणनाचे निर्गत परत करते.
\end{explain}

संगणनाचे प्रतिमान आंशिक पुनरावर्ती फलने हे असेल (क्लिनी यांनी
मूळतः हेच प्रतिमान वापरले होते), तर या प्रमेयावरून कोणतेही फलन
परिभाषित करण्यासाठी अपरिबद्ध शोधाचा, म्हणजे एकाच
$\umin{y}{f(\vec x, y)}$ प्रवर्तकाचा, एकदाच वापर पुरेसा असतो.
या अर्थाने प्रत्येक आंशिक पुनरावर्ती फलनाला प्रमाण रूप असते.

$T$ आणि $U$ वापरून $\cfind{0}$, $\cfind{1}$, $\cfind{2}$, \dots
हे प्रगणन परिभाषित करता येते. यापुढे $T$ आणि~$U$ यांची एक योग्य
निवड निश्चित केली आहे असे समजू, आणि
\[
\cfind{e}(x) \simeq U(\umin{s}{T(e,x,s)})
\]
हे समीकरण $\cfind{e}$ ची \emph{व्याख्या} मानू.

आणखी एक उपयुक्त तथ्य पुढीलप्रमाणे आहे:

\begin{thm}
प्रत्येक आंशिक संगणनक्षम फलनाला अनंत निर्देशांक असतात.
\end{thm}

हेही सहज स्पष्ट होते. एखाद्या संगणनक्षम फलनाचे वर्णन दिले असता,
त्याच फलनाचे (समान आदान-निर्गत जोड्यांचे) संगणन करणारे वेगळे वर्णन
तयार करता येते. उदाहरणार्थ, संगणनावर काहीही परिणाम न करणारी क्रिया
आधी करता येते ($0 = 0$ आहे का ते तपासणे किंवा $5$ पर्यंत मोजणे इत्यादी).
बदललेल्या वर्णनाचा निर्देशांक मूळ निर्देशांकापेक्षा नेहमी वेगळा असेल.
दोन्ही त्याच फलनाचे निर्देशांक आहेत; फक्त संगणनाची पद्धत किंचित वेगळी आहे.

\end{document}
